\section{Introduction}
Modern industrial infrastructures rely on rotating machinery, including wind turbines, aero-engines, and high-speed railways, where critical components such as bearings and gearboxes operate under demanding conditions \cite{zhai2021multi}. Reliable health monitoring is essential for operational safety and economic stability, as unexpected failures may lead to catastrophic accidents and significant financial losses \cite{lei2020applications}. Intelligent Fault Diagnosis (IFD) has emerged as a fundamental technology in smart manufacturing, using sensor-based monitoring to identify anomalies before they trigger system-wide failures. As industrial systems become increasingly complex, there is a growing demand for robust, autonomous, and high-precision diagnostic frameworks to ensure the continuity of critical production lines \cite{li2020deep}. Fig.~\ref{fig:concept} illustrates the proposed cross-modality semantic alignment approach for industrial zero-shot fault diagnosis.

Despite advancements in deep learning for fault diagnosis, most current methods depend on a closed-set assumption. This paradigm requires intensive manual labeling and a balanced distribution of training samples across all predefined health states \cite{zhang2022digital}. In practical industrial settings, machinery remains in normal states for the majority of its service life; consequently, fault samples---particularly those representing early-stage degradation---are rare and difficult to collect \cite{liu2018zero}. This imbalance creates a significant long-tail distribution problem where traditional supervised models often fail to generalize to unseen fault categories. To address data scarcity and the impracticality of exhaustive labeling, research has shifted from standard supervised learning toward Zero-Shot Learning (ZSL). ZSL aims to recognize novel fault classes without physical training samples by using auxiliary information to bridge known and unknown categories \cite{wang2019zero, socher2013zero}.

Current ZSL approaches in fault diagnosis primarily use semantic embeddings, such as manual attributes or word vectors (e.g., Word2Vec or GloVe), to associate vibration signals with class labels \cite{lampert2014attribute}. These methods face three critical limitations. First, defining manual attributes is labor-intensive and subjective, requiring experts to detail physical characteristics for every potential fault type \cite{fu2018zero}. Second, generic word embeddings are typically trained on web-scale natural language corpora, such as Wikipedia, which lack the technical specificity required for industrial diagnostics. For instance, the term ``outer race fault'' carries specific engineering connotations that generic vectors cannot capture with sufficient precision \cite{mikolov2013efficient}. Third, existing models often overlook the complex temporal dependencies in high-frequency vibration data, resulting in a domain gap where the alignment between signal features and abstract semantic space remains unstable \cite{xian2018zero, buchert2020zero}.

To address these challenges, this paper introduces the Cross-Modality Semantic Alignment Transformer (CMSA-Trans), a zero-shot fault diagnosis framework that uses Large Language Models (LLMs) to bridge the semantic divide. Unlike traditional methods, CMSA-Trans employs the descriptive capabilities of LLMs to generate high-fidelity, domain-specific fault semantics, transforming abstract class labels into detailed technical descriptions. By treating fault diagnosis as a cross-modal alignment task, the framework incorporates a Transformer-based temporal encoder to extract discriminative features from raw vibration signals. These features are projected into a shared latent space alongside the LLM-generated semantics. This alignment enables the model to learn the underlying physical relationships between signal patterns and their linguistic definitions, allowing for the detection of unseen faults based on the semantic proximity of the signal to known physical phenomena \cite{vaswani2017attention, radford2021learning}.

The primary contributions of this research are as follows. First, we develop the CMSA-Trans framework, which reformulates zero-shot fault diagnosis as a cross-modal semantic alignment task, eliminating the need for labeled samples of target fault categories. Second, we propose a systematic method to utilize LLMs for generating expert-level fault semantics, capturing the detailed physical characteristics of rotating machinery faults more effectively than traditional word embeddings. Third, a Transformer-based temporal feature extractor is designed to capture long-range dependencies in vibration signals, providing a robust representation for alignment with high-dimensional semantic descriptors. Finally, extensive experiments on multiple open-source and industrial datasets demonstrate that CMSA-Trans outperforms state-of-the-art ZSL methods, achieving high accuracy on unseen fault classes and showing strong generalization across various operating conditions.

The integration of LLMs, such as GPT-4, represents a shift in semantic feature engineering for industrial applications \cite{bubeck2023sparks}. Unlike static word embeddings that aggregate broad linguistic associations, LLMs can synthesize technical descriptions that capture fault physics, such as impact repetition frequencies and localized high-frequency resonances \cite{brown2020language}. This approach creates a richer and more discriminative semantic space, where the distance between fault classes reflects actual physical mechanisms rather than simple linguistic co-occurrence in general text. By utilizing these LLM-generated semantics, CMSA-Trans offers a data-efficient alternative to traditional ZSL approaches.

Robust cross-modal alignment is essential. In traditional deep learning models, features from sensor data are typically mapped directly to one-hot encoded labels, which lack inherent physical information and cannot generalize to unseen fault types \cite{goodfellow2016deep, lecun2015deep}. Our semantic alignment approach uses a projection layer to map the high-level signal representations from the Transformer into the same dimensional space as the LLM-derived technical descriptors. This ensures the model associates specific signal patterns, such as periodic transients or increases in broadband noise, with physical descriptions. When a new fault type with a unique technical signature occurs, the model can identify it by matching the signal pattern to its linguistic definition, even without labeled training examples.

Recent advancements in attention-based architectures have improved the processing of sequential data in various fields \cite{devlin2018bert, dosovitskiy2020image}. For rotating machinery, vibration signals are time-variant and exhibit long-range dependencies, where early-stage fault signatures may be obscured by noise or periodic components from normal operation \cite{zhang2021attention}. While Recurrent Neural Networks (RNNs) and Long Short-Term Memory (LSTM) units have been used to handle such data, they often face issues with vanishing gradients or capturing dependencies over the long windows typical of high-frequency samples \cite{hochreiter1997long}. The self-attention mechanism of the Transformer provides an effective method to weight different time steps according to their relevance, allowing the CMSA-Trans model to focus on the most informative transients and spectral characteristics across the signal window.

The use of large-scale pre-trained models for industrial semantic understanding is motivated by the extensive knowledge within these models \cite{zhao2023brain}. While an engineer might define a bearing fault by physical dimensions, an LLM can provide a textual profile that links those dimensions to expected sensor responses and characteristic frequencies. By aligning these descriptions with raw data, our approach moves closer to an expert-informed diagnostic system. This capability is vital for critical infrastructure where failures are rare but significant, as the ability to identify a new failure mode upon its first occurrence is essential for maintaining structural integrity.

Expanding zero-shot learning to include cross-modality also addresses the hubness problem often encountered in semantic embedding spaces, where certain vectors become nearest neighbors to many unrelated items in high-dimensional space \cite{radovanovic2010hubs}. A more descriptive and technically grounded semantic space provided by LLMs makes the separation between distinct fault modes---such as inner-race versus outer-race bearing defects---more pronounced. This enhanced class separation in the semantic domain improves diagnostic reliability and confidence scores for unseen categories. Validation on various machinery types, including complex planetary gearsets and varying load-speed conditions, confirms that the CMSA-Trans framework provides a scalable and robust solution for intelligent PHM (Prognostics and Health Management) systems.

The remainder of this paper is organized as follows. Section II reviews the related literature on intelligent fault diagnosis and zero-shot learning \cite{li2022zero}. Section III describes the CMSA-Trans architecture, including the LLM semantic generation and the cross-modal alignment mechanism. Section IV presents the experimental setup, results, and comparative analysis. Finally, Section V concludes the paper and discusses future research directions.
